\section*{Bab \ref{ch:b3:group-theory-applied} --- Penerapan Teori Grup}

\begin{solution}{exo:b3:group-theory-applied:1}
$n_{A_1} = \frac14(5 - 1 + 3 + 1) = 2$, $n_{A_2} = \frac14(5 - 1 - 3 - 1) = 0$,
$n_{B_1} = \frac14(5 + 1 + 3 - 1) = 2$, $n_{B_2} = \frac14(5 + 1 - 3 + 1) = 1$:
$\Gamma = 2A_1 + 2B_1 + B_2$, berdimensi 5.
\end{solution}

\begin{solution}{exo:b3:group-theory-applied:2}
Hitungan atom yang tetap di tempat sama dengan air: $\Gamma_{\mathrm{vib}} = 2A_1 + B_1$.
$A_1$ ($z$) dan $B_1$ ($x$) \omterm{def:b3:group-theory-applied:activity}{aktif IR}, dan keduanya juga \omterm{def:b3:group-theory-applied:activity}{aktif Raman} ($x^2$,
$xz$): tiga pita dalam setiap spektrum.
\end{solution}

\begin{solution}{exo:b3:group-theory-applied:3}
$B_1\otimes B_2$: \omterm{def:b3:point-groups:representation}{karakter} $(1, 1, -1, -1)$ $= A_2$. $A_2\otimes B_1$: $(1, -1, -1,
1) = B_2$. Dalam $C_{3v}$, $E\otimes E$ mempunyai \omterm{def:b3:point-groups:representation}{karakter} $(4, 1, 0)$;
reduksinya: $A_1 + A_2 +
E$.
\end{solution}

\begin{solution}{exo:b3:group-theory-applied:4}
Inframerah: kedua modus $t_2$, 3019 dan \qty{1306}{cm^{-1}}. Raman: keempatnya
($a_1$, $e$, $t_2$ semuanya bertransformasi seperti fungsi kuadratik). Tidak
ada \omterm{def:b3:point-groups:inversion}{pusat inversi}, sehingga kebetulan pita yang berimpit diizinkan.
\end{solution}

\begin{solution}{exo:b3:group-theory-applied:5}
Atom yang tetap di tempat $(4, 1, 2, 4, 1, 2)$ dikalikan
$\chi_{xyz} = (3, 0, -1, 1, -2, 1)$ memberikan
$\chi_{3N} = (12, 0, -2, 4, -2, 2)$, yang tereduksi menjadi $A_1' + A_2' + 3E' + 2A_2'' +
E''$. Kurangkan $E' + A_2''$ (translasi) dan $A_2' + E''$ (rotasi): $A_1' + 2E'
+ A_2''$. IR: $2E' + A_2''$, tiga pita; Raman: $A_1' + 2E'$, tiga pita. $A_2''$
(tekuk keluar bidang) hanya IR, $A_1'$ (ulur simetris) hanya Raman, dan
modus-modus $E'$ muncul dalam keduanya.
\end{solution}

\begin{solution}{exo:b3:group-theory-applied:6}
IR: kedua modus $T_{1u}$, dua pita. Raman: $A_{1g}$, $E_g$, $T_{2g}$, tiga pita.
$T_{2u}$ senyap. \ce{SF6} mempunyai \omterm{def:b3:point-groups:inversion}{pusat inversi}: eksklusi timbal balik.
\end{solution}

\begin{solution}{exo:b3:group-theory-applied:7}
Oleh $E$, $C_3$, $C_3^2$ fungsi $h_2$ dibawa ke $h_2$, $h_3$, $h_1$ (\omterm{def:b3:point-groups:representation}{karakter}
2, $-1$, $-1$); pencerminan menyumbang 0. $\hat P_Eh_2 \propto 2h_2 - h_3 - h_1$.
Tumpang-tindihnya dengan $2h_1 - h_2 - h_3$ (tumpang-tindih antaratom
diabaikan) adalah $-3$, dan fungsi yang terakhir itu bernorma $\sqrt6$;
setelah proyeksinya dikurangkan, $(2h_2 - h_1 - h_3) +
\frac12(2h_1 - h_2 - h_3) = \frac32(h_2 - h_3)$.
\end{solution}

\begin{solution}{exo:b3:group-theory-applied:8}
fac ($C_{3v}$, \omterm{def:b3:point-groups:representation}{karakter} $3, 0, 1$): $A_1 + E$, dua pita IR. mer ($C_{2v}$,
\omterm{def:b3:point-groups:representation}{karakter} $3, 1, 3, 1$): $2A_1 + B_1$, tiga pita. \ce{M(CO)6} ($O_h$): $A_{1g} + E_g
+ T_{1u}$, satu pita IR ($T_{1u}$).
\end{solution}

\begin{solution}{exo:b3:group-theory-applied:9}
\omterm{def:b3:point-groups:representation}{Karakter} $6, 0, 0, 2, 2, 0, 0, 0, 4, 2$: hanya $E$, $C_4$ dan $C_2$ sepanjang
sumbu-sumbu (dua ligan tetap di tempat), ketiga $\sigma_h$ (empat) dan keenam
$\sigma_d$ (dua) yang membiarkan ligan tetap di tempat. Reduksinya memberikan
$A_{1g} + E_g + T_{1u}$. Orbital $s$ logam ($a_{1g}$), $d_{z^2}$ dan $d_{x^2-y^2}$
($e_g$), serta $p_x$, $p_y$, $p_z$ ($t_{1u}$) berkombinasi; $t_{2g}$ tetap
nonikatan.
\end{solution}

\begin{solution}{exo:b3:group-theory-applied:10}
$n_i = \frac1h(1\cdot h\cdot\chi_i(E)) = d_i$. Dimensi representasi reguler
adalah $h$, dan juga $\sum_in_id_i = \sum_id_i^2$.
\end{solution}

\begin{solution}{exo:b3:group-theory-applied:11}
Dalam $D_{2h}$, dengan molekul terletak sepanjang sumbu $z$, banyaknya atom
yang tetap di tempat adalah $(4, 4, 0, 0, 0, 0, 4, 4)$, $\chi_{3N}
= (12, -4, 0, 0, 0, 0, 4, 4)$, yang tereduksi menjadi $2A_g + 2B_{2g} + 2B_{3g} + 2B_{1u} + 2B_{2u}
+ 2B_{3u}$. Kurangkan translasi $B_{1u} + B_{2u} + B_{3u}$ dan kedua rotasi
$B_{2g} + B_{3g}$: $2A_g + B_{2g} + B_{3g} + B_{1u} + B_{2u} + B_{3u}$, tujuh modus.
Dalam $D_{\infty h}$: $2\Sigma_g^+ + \Pi_g + \Sigma_u^+ + \Pi_u$. IR: $\Sigma_u^+$ dan
$\Pi_u$ (dua pita); Raman: $2\Sigma_g^+$ dan $\Pi_g$ (tiga pita); tidak ada
yang berimpit.
\end{solution}

\begin{solution}{exo:b3:group-theory-applied:12}
$1b_2 \to 4a_1$: $B_2$, $y$, diizinkan, tegak lurus pada bidang. $3a_1 \to 4a_1$:
$A_1$, $z$, diizinkan, sepanjang sumbu $C_2$. $1b_2 \to 2b_1$: $A_2$, terlarang.
$1b_1 \to 4a_1$: $B_1$, $x$, diizinkan, pada bidang molekul dan tegak lurus
pada sumbu.
\end{solution}

\begin{solution}{pb:b3:group-theory-applied:1}
\textbf{1.} cis: kedua CO pada titik sudut yang bersebelahan; trans: yang
berhadapan; fac: ketiga CO pada satu sisi; mer: ketiga CO pada satu meridian
(dua saling trans, satu di antaranya).
\textbf{2.} $C_{2v}$.
\textbf{3.} $D_{4h}$.
\textbf{4.} $C_{3v}$.
\textbf{5.} $C_{2v}$.
\textbf{6.} Hanya trans-\ce{ML4(CO)2}.
\textbf{7.} Operasi yang memindahkan sebuah CO ke CO lain menempatkan
vektornya di luar diagonal (sumbangan 0); operasi yang membiarkan sebuah CO
di tempat memetakan vektor ikatannya ke dirinya sendiri (+1); tidak ada
operasi yang membalik vektor C--O.
\textbf{8.} \omterm{def:b3:point-groups:representation}{Karakter} $(2, 0, 2, 0)$, dengan $\sigma_v(xz)$ bidang yang memuat kedua
CO: $A_1 +
B_1$.
\textbf{9.} \omterm{def:b3:point-groups:representation}{Karakter} 2 untuk $E$, $C_4$, $C_2$, $\sigma_v$, $\sigma_d$, dan 0 untuk
yang lain: $A_{1g} + A_{2u}$.
\textbf{10.} $(3, 0, 1)$: $A_1 + E$.
\textbf{11.} $(3, 1, 3, 1)$: $2A_1 + B_1$.
\textbf{12.} 2, 2, 3, 3: satu dimensi untuk setiap CO.
\textbf{13.} cis 2, trans 1, fac 2, mer 3.
\textbf{14.} cis 2, trans 1, fac 2, mer 3.
\textbf{15.} trans: pita IR ($A_{2u}$) dan pita Raman ($A_{1g}$) isomer itu
berasal dari modus yang berbeda.
\textbf{16.} $A_{2u}$ berganti tanda oleh $i$: satu CO memanjang sementara yang
lain memendek, berlawanan fase.
\textbf{17.} $E$ adalah pasangan terdegenerasi: dua modus berfrekuensi sama
memberikan satu pita, ditambah pita $A_1$.
\textbf{18.} Isomer tersubstitusi-trans yang sentrosimetris.
\textbf{19.} Tiga pita IR: mer (fac akan memberikan dua).
\textbf{20.} Ya: mer mempunyai tiga modus CO yang \omterm{def:b3:group-theory-applied:activity}{aktif Raman} pada frekuensi
yang sama, karena tidak sentrosimetris; fac akan memperlihatkan dua.
\textbf{21.} $A_1$: vibrasi ulur sefase bersifat simetris total.
\textbf{22.} Campuran fac/mer akan memperlihatkan sampai lima pita IR, dengan
pita-pita fac pada posisi yang berbeda; tiga pita yang bersih cocok dengan
satu isomer saja. Spektrum NMR ${}^{13}$C akan memutuskannya: dua sinyal
karbonil dengan rasio 2:1 untuk mer, satu sinyal untuk fac.
\textbf{23.} \textbf{Isomer mer mempunyai tiga vibrasi ulur CO yang \omterm{def:b3:group-theory-applied:activity}{aktif IR},
isomer fac dua}: produknya adalah mer.
\end{solution}
